\begin{table}[htbp]
\centering
\caption{Top 10 Departamentos: Mayor Cambio en Defección de Cabildo Abierto (2019--2024)}
\label{tab:temporal_departments_top10}
\small
\begin{tabular}{lccc}
\toprule
{\textbf{Departamento}} & {\textbf{Def. 2019 (\%)}} & {\textbf{Def. 2024 (\%)}} & {\textbf{Cambio (pp)}} \\
\midrule
1. Colonia & 30.2 & 52.6 & +22.3 \\
 & $[26.5, 33.9]$ & $[29.4, 77.0]$ & {\footnotesize Parcial} \\
2. Treinta y Tres & 58.4 & 70.2 & +11.9 \\
 & $[54.5, 62.2]$ & $[39.1, 94.5]$ & {\footnotesize Parcial} \\
3. Rocha & 36.7 & 40.0 & +3.4 \\
 & $[35.4, 38.1]$ & $[9.0, 71.7]$ & {\footnotesize Parcial} \\
4. Artigas & 34.0 & 34.4 & +0.3 \\
 & $[29.9, 38.2]$ & $[19.5, 48.9]$ & {\footnotesize Parcial} \\
5. Río Negro & 96.6 & 95.7 & -0.9 \\
 & $[94.7, 98.2]$ & $[87.8, 99.5]$ & {\footnotesize Parcial} \\
6. Flores & 27.9 & 23.2 & -4.7 \\
 & $[18.6, 38.1]$ & $[1.0, 68.9]$ & {\footnotesize Parcial} \\
7. Rivera & 46.0 & 24.4 & -21.6 \\
 & $[43.1, 49.0]$ & $[11.2, 39.1]$ & {\footnotesize Sí} \\
8. Cerro Largo & 42.6 & 14.8 & -27.8 \\
 & $[40.4, 44.8]$ & $[2.0, 29.4]$ & {\footnotesize Sí} \\
9. Maldonado & 78.0 & 47.9 & -30.1 \\
 & $[73.7, 82.5]$ & $[21.0, 75.6]$ & {\footnotesize Parcial} \\
10. Florida & 83.6 & 38.7 & -44.9 \\
 & $[79.3, 87.9]$ & $[0.9, 74.9]$ & {\footnotesize Sí} \\
\bottomrule
\end{tabular}
\begin{flushleft}
\small \textit{Nota:} Departamentos ordenados por mayor aumento en defección de Cabildo Abierto hacia Frente Amplio. Segunda fila muestra intervalos creíbles al 95\%. Significancia indica si intervalos no se solapan.
\end{flushleft}
\end{table}
